\chapter{Conclusão}
\label{cap:05}

O discurso de ódio disseminado por meio de memes representa um dos desafios mais complexos da moderação de conteúdo digital. Diferentemente de textos ou imagens isolados, os memes constroem seu significado na justaposição intencional entre as duas modalidades, e é exatamente essa característica que torna as abordagens unimodais insuficientes e que motivou a proposta desta pesquisa.

Esta pesquisa investigou a identificação automatizada de discurso de ódio em memes por meio de uma arquitetura progressiva de três experimentos, cada um construído sobre as limitações do anterior. O Experimento~1 estabeleceu o modelo CLIP ViT-L/14 em modo \textit{zero-shot}, embora seja um dos modelos de alinhamento multimodal mais avançados disponíveis, opera próximo ao acaso nessa tarefa (AUROC 0,5392), demonstrando que a similaridade semântica geral não é suficiente para detectar a incongruência semântica específica dos memes de ódio. O Experimento~2 mostrou que a introdução de supervisão, ainda com o CLIP congelado e enriquecido por índices de agrupamento semântico, foi o fator determinante para o salto de desempenho (AUROC 0,7318), com a contribuição dos \textit{clusters} sendo incremental sobre a base multimodal supervisionada. O Experimento~3 avançou ao modelar explicitamente as correlações entre as dimensões visuais e textuais por meio da \textit{Feature Interaction Matrix}, alcançando AUROC de 0,8121.

A lógica \textit{fuzzy}, integrada em todos os experimentos e evoluindo de um classificador baseado em entropia de Shannon para um sistema Mamdani com antecedentes visuais e textuais independentes, demonstrou que a gradação linguística do ódio é computacionalmente viável. Mais do que melhorar métricas, o \textit{fuzzy} qualifica a decisão: ao distribuir casos de alta incerteza em categorias como \textit{Ambiguidade/Moderado}, o sistema sinaliza quais memes demandam revisão humana em vez de forçar uma classificação binária sobre conteúdo inerentemente ambíguo.

Retomando o objetivo geral desta pesquisa, os resultados demonstram que a fusão cruzada supervisionada, com representações normalizadas por modalidade e complementada por interpretabilidade \textit{fuzzy}, produz um sistema competitivo com a literatura (AUROC de 0,8121, frente a 0,8265 de desempenho humano e 0,858 do Hate-CLIPper original), utilizando uma arquitetura replicável e recursos computacionais acessíveis. Além dos resultados experimentais, este trabalho produziu um artefato prático utilizável: o modelo do Experimento~3 foi disponibilizado publicamente no \textit{Hugging Face Spaces}, com interface interativa que permite a qualquer usuário submeter um meme e obter tanto a classificação binária quanto a graduação \textit{fuzzy} do grau de ódio, demonstrando a viabilidade prática da abordagem para além do ambiente experimental do Google Colab. O código, os \textit{notebooks} e os resultados estão documentados de forma reproduzível, o que permite que trabalhos futuros partam diretamente dos \textit{checkpoints} treinados.

A Tabela~\ref{tab:comparacao-literatura} compara o AUROC do modelo proposto com resultados reportados na literatura.

\begin{table}[htbp]
\centering
\caption{Comparação do resultado final com a literatura.}
\label{tab:comparacao-literatura}
\begin{tabular}{lc}
\hline
\textbf{Referência} & \textbf{AUROC} \\
\hline
Desempenho humano \cite{metaai2020} & 0,8265 \\
VisualBERT \cite{ArticleHatefulMemesChallenge2020} & 0,730 \\
HateDetectron \cite{velioglu2020} & 0,811 \\
\textbf{Este trabalho (Experimento 3)} & \textbf{0,8121} \\
Hate-CLIPper original \cite{kumar2022} & 0,858 \\
\hline
\end{tabular}
\end{table}

Os desafios que permanecem em aberto delineiam três caminhos naturais para pesquisas futuras:

\begin{itemize}
    \item construção de \textit{datasets} multilíngues e multiculturais, especialmente para o contexto brasileiro;
    \item exploração de modelos de linguagem multimodais de grande porte (GPT-4V, LLaVA, Gemini) para lidar com sarcasmo e ganhar interpretabilidade;
    \item incorporação de técnicas de \textit{debiasing} e auditoria de equidade por subgrupo demográfico.
\end{itemize}

O \textit{Hateful Memes Dataset} é composto exclusivamente por memes em inglês com referências culturais ocidentais, o que limita a generalização para contextos como o brasileiro, onde gírias regionais, ironia contextual e referências políticas locais alteram substancialmente o significado do conteúdo. A construção de \textit{datasets} multilíngues e multiculturais anotados é um passo necessário para tornar sistemas como este viáveis em escala global.

Nos Experimentos~1 e~2, a baixa revocação para a classe de ódio evidenciou que o modelo deixava de detectar a maioria dos casos genuinamente odiosos. O Experimento~3 inverteu esse padrão ao otimizar o limiar para maximizar a revocação (0,86), às custas de uma precisão mais moderada (0,55), uma escolha adequada para cenários de triagem, mas que ainda reflete uma limitação comum a todos os experimentos: mesmo quando detecta o conteúdo odioso, o modelo não captura adequadamente a justaposição semântica contraditória entre imagem e texto por meio de raciocínio explícito, apoiando-se em correlações estatísticas aprendidas em vez de compreensão contextual genuína. Modelos de Linguagem de Grande Porte multimodais (GPT-4V, LLaVA, Gemini) representam a próxima fronteira natural, pois combinam compreensão profunda de contexto com capacidade de explicar a classificação, o que endereçaria simultaneamente os problemas de sarcasmo e interpretabilidade.

Um risco central em sistemas de moderação automatizada é a reprodução e amplificação de preconceitos presentes nos dados de treinamento. No contexto de detecção de ódio, \textit{datasets} desbalanceados podem levar o modelo a associar características identitárias a conteúdo ofensivo, gerando falsos positivos desproporcionais para grupos historicamente marginalizados. Trabalhos futuros podem incorporar técnicas de \textit{debiasing}, auditorias de equidade por subgrupo e avaliação de calibração por classe demográfica.

A automação da moderação levanta dilemas que vão além da acurácia: privacidade na análise em escala de conteúdo pessoal, risco de censura de expressões legítimas e a questão de quem define os limites do ódio. A interpretabilidade fornecida pela lógica \textit{fuzzy} é um passo na direção da responsabilização, mas não substitui a supervisão humana, especialmente para os casos enquadrados como \textit{Ambiguidade/Moderado}, que por definição demandam julgamento contextual. Um sistema de moderação responsável é, necessariamente, um sistema híbrido.
